\subsection{Cryptographic Algorithms}\label{sec:cryptography_basics:algorithms}

A \textbf{cryptographic algorithm} (or simply algorithm) is a function --- that is, having one or more inputs and returning one or more outputs --- used in the context of cryptography. Many algorithms rely on one-way and trapdoor functions (\Cref{sec:trapdoor_functions}). An \textbf{encryption algorithm} converts \textbf{plaintexts} --- that is, intelligible data --- into \textbf{ciphertexts} --- that is, unintelligible data. A ciphertext may also be called \textit{encrypted plaintext}. A plaintext is denoted with $m$, and the set of all plaintexts with $\Plaintexts$. Similarly, a ciphertext is denoted with $c$, and the set of all ciphertexts with $\Ciphertexts$. An encryption algorithm is (practically) always paired with a \textbf{decryption algorithm} which converts ciphertexts back into the corresponding plaintexts. A pair of encryption and decryption algorithms has the \textbf{correctness} property if decrypting any ciphertext using correct and expected parameters (e.g., keys) returns the corresponding plaintext. Ideally, any pair of encryption and decryption algorithms possesses the correctness property. Finally, there exist also \textbf{translation algorithms} which translate ciphertexts --- that is, plaintexts encrypted under a specific key --- into different ciphertexts --- that is, the same plaintexts encrypted under a different key; translation does not imply decryption and can be seen as a particular instance of re-encryption (\Cref{sec:TEMPORARY_LABEL:re_encryption}).